\documentclass[10pt, letterpaper]{article}

% Packages:
\usepackage[
    ignoreheadfoot, % set margins without considering header and footer
    top=2 cm, % seperation between body and page edge from the top
    bottom=2 cm, % seperation between body and page edge from the bottom
    left=2 cm, % seperation between body and page edge from the left
    right=2 cm, % seperation between body and page edge from the right
    footskip=1.0 cm, % seperation between body and footer
    % showframe % for debugging 
]{geometry} % for adjusting page geometry
\usepackage[explicit]{titlesec} % for customizing section titles
\usepackage{tabularx} % for making tables with fixed width columns
\usepackage{array} % tabularx requires this
\usepackage[dvipsnames]{xcolor} % for coloring text
\definecolor{primaryColor}{RGB}{0, 79, 144} % define primary color
\usepackage{enumitem} % for customizing lists
\usepackage{fontawesome5} % for using icons
\usepackage{amsmath} % for math
\usepackage[
    pdftitle={Bettina Ballin's CV},
    pdfauthor={Bettina Ballin},
    pdfcreator={LaTeX with RenderCV},
    colorlinks=true,
    urlcolor=primaryColor
]{hyperref} % for links, metadata and bookmarks
\usepackage[pscoord]{eso-pic} % for floating text on the page
\usepackage{calc} % for calculating lengths
\usepackage{bookmark} % for bookmarks
\usepackage{wrapfig} % for wrapfigure environment
\usepackage{lastpage} % for getting the total number of pages
\usepackage{changepage} % for one column entries (adjustwidth environment)
\usepackage{paracol} % for two and three column entries
\usepackage{ifthen} % for conditional statements
\usepackage{needspace} % for avoiding page brake right after the section title
\usepackage{iftex} % check if engine is pdflatex, xetex or luatex
\usepackage{graphicx}
\usepackage{tikz}
% Ensure that generate pdf is machine readable/ATS parsable:
\ifPDFTeX
    \input{glyphtounicode}
    \pdfgentounicode=1
    \usepackage[T1]{fontenc}
    \usepackage[utf8]{inputenc}
    \usepackage{lmodern}
\fi

\usepackage[default, type1]{sourcesanspro} 

% Some settings:
\AtBeginEnvironment{adjustwidth}{\partopsep0pt} % remove space before adjustwidth environment
\pagestyle{empty} % no header or footer
\setcounter{secnumdepth}{0} % no section numbering
\setlength{\parindent}{0pt} % no indentation
\setlength{\topskip}{0pt} % no top skip
\setlength{\columnsep}{0.15cm} % set column seperation
\makeatletter
\let\ps@customFooterStyle\ps@plain % Copy the plain style to customFooterStyle
\patchcmd{\ps@customFooterStyle}{\thepage}{
    \color{gray}\textit{\small Bettina Ballin - Seite \thepage{} von \pageref*{LastPage}}
}{}{} % replace number by desired string
\makeatother
\pagestyle{customFooterStyle}

\titleformat{\section}{
    % avoid page braking right after the section title
    \needspace{4\baselineskip}
    % make the font size of the section title large and color it with the primary color
    \Large\color{primaryColor}
}{
}{
}{
    % print bold title, give 0.15 cm space and draw a line of 0.8 pt thickness
    % from the end of the title to the end of the body
    \textbf{#1}\hspace{0.15cm}\titlerule[0.8pt]\hspace{-0.1cm}
}[] % section title formatting

\titlespacing{\section}{
    % left space:
    -1pt
}{
    % top space:
    0.3 cm
}{
    % bottom space:
    0.2 cm
} % section title spacing

% \renewcommand\labelitemi{$\vcenter{\hbox{\small$\bullet$}}$} % custom bullet points
\newenvironment{highlights}{
    \begin{itemize}[
        topsep=0.10 cm,
        parsep=0.10 cm,
        partopsep=0pt,
        itemsep=0pt,
        leftmargin=0.4 cm + 10pt
    ]
}{
    \end{itemize}
} % new environment for highlights

\newenvironment{highlightsforbulletentries}{
    \begin{itemize}[
        topsep=0.10 cm,
        parsep=0.10 cm,
        partopsep=0pt,
        itemsep=0pt,
        leftmargin=10pt
    ]
}{
    \end{itemize}
} % new environment for highlights for bullet entries


\newenvironment{onecolentry}{
    \begin{adjustwidth}{
        0.2 cm + 0.00001 cm
    }{
        0.2 cm + 0.00001 cm
    }
}{
    \end{adjustwidth}
} % new environment for one column entries

\newenvironment{twocolentry}[2][]{
    \onecolentry
    \def\secondColumn{#2}
    \setcolumnwidth{\fill, 4.5 cm}
    \begin{paracol}{2}
}{
    \switchcolumn \raggedleft \secondColumn
    \end{paracol}
    \endonecolentry
} % new environment for two column entries

\newenvironment{threecolentry}[3][]{
    \onecolentry
    \def\thirdColumn{#3}
    \setcolumnwidth{1 cm, \fill, 4.5 cm}
    \begin{paracol}{3}
    {\raggedright #2} \switchcolumn
}{
    \switchcolumn \raggedleft \thirdColumn
    \end{paracol}
    \endonecolentry
} % new environment for three column entries

\newenvironment{header}{
    \setlength{\topsep}{0pt}\par\kern\topsep\centering\color{primaryColor}\linespread{1.5}
}{
    \par\kern\topsep
} % new environment for the header

\newcommand{\placelastupdatedtext}{% \placetextbox{<horizontal pos>}{<vertical pos>}{<stuff>}
  \AddToShipoutPictureFG*{% Add <stuff> to current page foreground
    \put(
        \LenToUnit{\paperwidth-2 cm-0.2 cm+0.05cm},
        \LenToUnit{\paperheight-1.0 cm}
    ){\vtop{{\null}\makebox[0pt][c]{
        \small\color{gray}\textit{Zuletzt aktualisiert: Juli 2025}\hspace{\widthof{Zuletzt aktualisiert: Juli 2025}}
    }}}%
  }%
}%

% save the original href command in a new command:
\let\hrefWithoutArrow\href

% new command for external links:
\renewcommand{\href}[2]{\hrefWithoutArrow{#1}{\ifthenelse{\equal{#2}{}}{ }{#2 }\raisebox{.15ex}{\footnotesize \faExternalLink*}}}


\begin{document}
    \newcommand{\AND}{\unskip
        \cleaders\copy\ANDbox\hskip\wd\ANDbox
        \ignorespaces
    }
    \newsavebox\ANDbox
    \sbox\ANDbox{}

    \placelastupdatedtext
    \begin{header}

    \begin{wrapfigure}{R}{0.15\textwidth}
        \begin{center}
        \vspace{-0.8cm}
        \hspace{-9cm}
        \begin{tikzpicture}
        \sbox0{\includegraphics[width=0.2\textwidth]{profile2.jpg}}%
        \path[clip,draw,rounded corners=0.5cm] (0,0) rectangle (\wd0,\ht0);
\path (0.5\wd0,0.5\ht0) node[inner sep=0pt]{\usebox0};
        \end{tikzpicture}
        \end{center}
    \end{wrapfigure}
    
        \fontsize{30 pt}{30 pt}
        \textbf{Bettina Ballin}
    
        \vspace{0.3 cm}

        \normalsize
        \mbox{{\footnotesize\faMapMarker*}\hspace*{0.13cm}Lotichiusstr. 43, 36381 Schlüchtern}%
        \kern 0.25 cm%
        \AND%
        \kern 0.25 cm%
        \mbox{\hrefWithoutArrow{mailto:betty.ballin@yahoo.de}{{\footnotesize\faEnvelope[regular]}\hspace*{0.13cm}betty.ballin@yahoo.de}}%
        \kern 0.25 cm%
        \AND%
        \kern 0.25 cm%
        \mbox{\hrefWithoutArrow{tel:+49-176-808-42-853}{{\footnotesize\faPhone*}\hspace*{0.13cm}+49 176 808 42 853}}%
        \kern 0.25 cm%
        \AND%
        \kern 0.25 cm%
        \mbox{\hrefWithoutArrow{https://bballin.me/}{{\footnotesize\faLink}\hspace*{0.13cm}bballin.me}}%
        \kern 0.25 cm%
        \AND%
        \kern 0.25 cm%
        \mbox{\hrefWithoutArrow{https://linkedin.com/in/bettina-ballin-792b5b167}{{\footnotesize\faLinkedinIn}\hspace*{0.13cm}bettina-ballin-792b5b167}}%
        \kern 0.25 cm%
        \AND%
        \kern 0.25 cm%
        \mbox{\hrefWithoutArrow{https://github.com/bettyballin}{{\footnotesize\faGithub}\hspace*{0.13cm}bettyballin}}%

         \vspace{0.5cm}
         
    \end{header}

    \vspace{0.3 cm - 0.3 cm}


\section{Berufserfahrung}
\begin{twocolentry}{
    Remote / Deutschland

    Apr 2025 – heute
}
    \textbf{Selbstständig/IT-Freelancer} im Bereich Software \& Consulting
    \begin{highlights}
        \item Architektur und Optimierung maßgeschneiderter Softwarelösungen für KMU
        \item Administration von Webprojekten und Hosting der IT-Infrastruktur
    \end{highlights}
\end{twocolentry}
\begin{twocolentry}{
    Darmstadt

    Nov 2022 – Apr 2025
}
    \textbf{Fraunhofer SIT}, Wiss. Mitarbeiter im Bereich Secure Software Engineering
    \begin{highlights}
        \item Veröffentlichungen für den Bereich KI und Cybersicherheit
        \item Forschung an Anomalie-Erkennungssystemen für deutsche Raumfahrtunternehmen
        \item Mitarbeit an Penetrationstests für staatliche Unternehmen sowie Erstellung von wissenschaftlichen Arbeiten für Forschungspartner
    \end{highlights}
\end{twocolentry}

\begin{twocolentry}{
    Darmstadt

    Feb 2018 – Jun 2020
}
    \textbf{Software AG}, webMethods Consultant
    \begin{highlights}
        \item IT-Berater in agilen IT-Teams führender deutscher Banken und Versicherungen
        \item Implementierung von Backend- und Frontend-Lösungen mithilfe der webMethods-Integrationsplattform
    \end{highlights}
\end{twocolentry}

\hspace{0.15cm}
\begin{twocolentry}{

    Mai 2022 – Sep 2022
}
    \textbf{Technische Universität Darmstadt}, Studentische Hilfskraft
    \begin{highlights}
        \item \textbf{Tutorium:} Statistical Machine Learning, Unterstützung des Übungsbetriebs
    \end{highlights}
\end{twocolentry}

\begin{twocolentry}{

    Okt 2020 – Apr 2022
}
    \begin{highlights}
        \item \textbf{PTW Institut}: Implementierung von Energieflexibilitäts-Tools in Python, Entwicklung von Webapplikationen mithilfe von Flask und Docker
    \end{highlights}
\end{twocolentry}


\begin{twocolentry}{

    Mär 2016 – Dez 2017
}
    \begin{highlights}
        \item \textbf{SECUSO Forschungsgruppe}: Entwicklung von Thunderbird-Add-Ons zum Schutz gegen Phishing-Angriffe
    \end{highlights}
\end{twocolentry}

\hspace{0.2cm}
\begin{twocolentry}{
    Rhein-Main-Gebiet, Deutschland

    Okt 2014 – Mär 2015
}
    \textbf{Staufenbiel Institut}, Event-Hostess \& Promotion
    \begin{highlights}
        \item Recruiting und Veranstaltungsorganisation mit Fokus auf junge Hochschulabsolventen
    \end{highlights}
\end{twocolentry}

\section{Ausbildung}

\begin{threecolentry}{\textbf{M.Sc.}}{
    Okt 2020 – heute
}
    \textbf{Technische Universität Darmstadt}
    \begin{highlights}
        \item M.Sc. Informatik: Note: 1,3 (sehr gut)
        \item Masterarbeit: \textit{Robustheit von Membership-Inference-Angriffen auf neuronale Netzwerke mit Label-Only-Information}
        \item Schwerpunkt auf KI-Technologien, Nebenfach: Philosophie
    \end{highlights}
\end{threecolentry}

\begin{threecolentry}{\textbf{B.Sc.}}{
    Okt 2013 – Dez 2017
}
    \textbf{Technische Universität Darmstadt}
    \begin{highlights}
        \item B.Sc. Informatik: Note: 1,9 (gut)
        \item Bachelorarbeit: \textit{Toolgestützte Erkennung schädlicher Links in Online-E-Mail-Applikationen}
    \end{highlights}
\end{threecolentry}

\begin{threecolentry}{\textbf{Abitur}}{
    Aug 2005 – Jul 2013
}
    \textbf{Ulrich-von-Hutten-Gymnasium, Schlüchtern}
    \begin{highlights}
        \item Allgemeine Hochschulreife
    \end{highlights}
\end{threecolentry}
\section{Zertifikate und Auszeichnungen}

\begin{highlights}
    \item T.I.S.P. TeleTrust Information Security Professional (2023)
    \item Deutschlandstipendium (2020)
    \item webMethods 9.9 Certified ESB Developer (2018)
\end{highlights}

\section{Kenntnisse und Fähigkeiten}

\begin{highlights}
    \item \textbf{Programmiersprachen:} Java, Python, C/C++, C\#, SQL, XML, LaTeX
    \item \textbf{Web-Entwicklung:} HTML, CSS, JavaScript, jQuery, Angular, Flask
    \item \textbf{Technologien:} Docker, Android Studio, Eclipse, Visual Studio Code, GitHub, JIRA, Confluence, Microsoft Office
    \item \textbf{Datenbanken:} Oracle, MongoDB
    \item \textbf{KI und Maschinelles Lernen:} Deep Learning, TensorFlow, PyTorch, NLP, Computer Vision, Anomalieerkennung
    \item \textbf{IT-Sicherheit:} Cybersecurity, Penetrationstests, Schwachstellenanalyse
    \item \textbf{Betriebssysteme:} Linux, Windows
\end{highlights}

\section{Sprachkenntnisse}

\begin{highlights}
    \item Deutsch: Muttersprache
    \item Englisch: Fließend
    \item Französisch, Spanisch: Erweiterte Grundkenntnisse
    \item Chinesisch: Grundkenntnisse
\end{highlights}

\section{Sonstiges}

\begin{highlights}
    \item Organisation und Mentoring in der Orientierungswoche für Erstsemester an der TU Darmstadt (Wintersemester 2016)
    \item Interessen: Klavier, Logikrätsel, Yoga, Bouldern
    \item Stärken: Schnelle Auffassungsgabe, starkes Interesse an mathematisch-technischen Herausforderungen, Teamfähigkeit und Kommunikation
\end{highlights}


\end{document}